\chapter{Aseguramiento de la Calidad y Validación}\label{calidad_pruebas}

\azul{Este capítulo consolida los planes y técnicas aplicados para verificar que el artefacto desarrollado cumpla con los estándares de calidad, seguridad y funcionamiento estipulados. La calidad en ingeniería se define como el grado en que un conjunto de características inherentes cumple con los requisitos especificados \cite[p.~14]{iso29148_2018},\cite[p.~412]{pressman2010ingenieria}.}

\section{Aseguramiento de la Calidad (QA) (\rojo{Opcional})}

\azul{Define el marco normativo y las técnicas utilizadas para garantizar el cumplimiento de los estándares de la disciplina.}

\subsection{Modelo o Estándar de Calidad}

\azul{Se especifica el modelo de referencia adoptado:}
\begin{itemize}
	\item \azul{\textbf{Exclusivo de Informática:} Modelos de calidad de software como ISO/IEC 25010 (SQuaRE) o McCall, evaluando atributos de mantenibilidad, usabilidad y eficiencia.}
	\item \azul{\textbf{Otras Ingenierías:} Estándares ISO 9001, normas chilenas de construcción (NCh), especificaciones ASTM/ASME o normativas de seguridad ambiental.}
\end{itemize}

\subsection{Técnica y Plan de Implementación}

\azul{Describe las auditorías, revisiones de código, inspecciones técnicas en terreno o controles de procesos que se ejecutaron de manera sistemática.}

\section{Proceso de Prueba y Ensayos de Verificación (\rojo{Opcional})}\label{prueba}

\azul{Detalla la estrategia experimental utilizada para constatar la ausencia de fallas en el diseño o implementación.}

\subsection{Criterios y Enfoques de Prueba}

\azul{Se definen las filosofías de evaluación aplicadas:}
\begin{itemize}
	\item \azul{\textbf{Enfoque de Caja Blanca (Estructural):} En \textbf{Informática}, consiste en la revisión directa de las líneas de código e itinerarios lógicos. En \textbf{otras ingenierías}, equivale al análisis de esfuerzo interno, inspección ultrasónica de soldaduras o ensayos no destructivos.}
	\item \azul{\textbf{Enfoque de Caja Negra (Funcional):} Evaluación de las respuestas o salidas del sistema ante entradas específicas, desconociendo la estructura interna.}
\end{itemize}

\section{Niveles de Prueba y Ensayos Sistemáticos}

\azul{Se detallan las pruebas ejecutadas en orden jerárquico de complejidad:}

\subsection{Pruebas Unitarias o Ensayos de Componentes}

\azul{\textbf{Exclusivo de Informática / Electrónica:} Verificación aislada de funciones, clases o módulos de código individual. En proyectos de infraestructura u obras físicas, esta subsección equivale a los ensayos de resistencia de probetas de hormigón o calibración individual de sensores.}

\subsection{Pruebas de Integración}

\azul{Evaluación de la interacción entre dos o más módulos, subsistemas o componentes acoplados. En software, se evalúa la comunicación entre la API y la base de datos; en otras ingenierías, la unión entre elementos estructurales o mecánicos.}

\subsection{Pruebas de Sistema, Rendimiento y Resistencia}

\azul{Evaluación global del artefacto en condiciones límite o bajo carga operacional:}
\begin{itemize}
	\item \azul{\textbf{Pruebas de Rendimiento y Carga:} Medición de tiempos de respuesta bajo peticiones masivas o evaluación de deformaciones bajo carga máxima.}
	\item \azul{\textbf{Pruebas de Seguridad y Resistencia:} Vulnerabilidades informáticas o tolerancia del material ante estrés ambiental/térmico.}
	\item \azul{\textbf{Pruebas de Recuperación:} Capacidad del sistema para restablecer su operación tras una falla imprevista.}
\end{itemize}

\subsection{Pruebas de Aceptación y Validación por el Usuario}

\azul{Instancia final en la que los usuarios, clientes o contrapartes técnicas evalúan el producto en un entorno real o de marcha blanca, validando que satisface la problemática identificada en el Capítulo \ref{formulacion} \cite[p.~240]{sommerville2011ingenieria}.}

\moradoc{Ejemplo: Se ejecutó una prueba de aceptación con tres operadores del comité de agua rural, registrando un 95\% de satisfacción en la usabilidad del panel de control y confirmando la correcta recepción de las alertas telemétricas.}